% Concrete application flows, checked against apps/fwdbase/fb_opaque.c,
% apps/obench/obench.c, apps/okvs/okvs.c and the evaluation plans.
% Implemented benchmark modes are not evidence for unmeasured application scenarios.
\section{Applications}
\label{sec:apps}

\paragraph{Named forwarding.} Our framed benchmark relay parses a copied length header,
discards it and captures the body. It sends the object to a configured destination with a
generated length prefix, marking the send as its last use. The parser uses the same stream
library with copying or object storage.

The forwarding experiments use a byte-stream relay: query available bytes, capture them
and send with the last-use flag, in both directions. A limit on bytes awaiting completion
lets a slow destination slow input through TCP flow control.

\paragraph{Mapped forwarding and processing.} The implemented bound relay uses the same
operations with a store bound to the I/O area, forwarding unchanged bytes. A trusted
application could query an object's extents, validate the area offsets and read those bytes
directly. The kernel reclaims buffers after all holders release them; no raw receive
references are returned by the application.

An editing extension, neither implemented nor measured, could replace a fixed-width
application field while the relay exclusively owns the captured range. Keeping the length
preserves framing. The relay would send by handle, avoid edits while those bytes might
still be in use, and stop accessing extents before releasing or transferring the handle.

The implemented raw mapped relay instead uses \recvzc and \sendzc, pins each send's buffers
and returns receive references after buffer-reuse notifications. Copying, \splice and sockmap
relays provide the comparisons in \S\ref{sec:eval:fwd}.

\paragraph{Named storage.} Our store implements the memcached text protocol. A
\code{SET} parses the command line to obtain the key and value length,
discards the command bytes, and captures the value as an object. It validates the trailing
delimiter before adding the value to the index. The application keeps the key, expiry time
and replacement policy; the kernel keeps the value's bytes.

On a \code{GET} hit, it sends the held value with a generated response header and trailer.
Repeated readers share the value. A pending reply keeps an application reference until its
send completes, so replacement or eviction does not invalidate that reply. Expired entries
are misses; eviction follows a least-recently-used policy adjusted under memory pressure.

The same server selects copying, named or mapped-bound storage while keeping its parser and
index. An optional small-value path keeps values in application memory. The public build has
one worker. An experimental multiworker version shares the keyspace and budget; its recorded
named/copy functional checks do not qualify performance, and mapped qualification remains pending.

\paragraph{Further workloads.} The benchmark tool implements retention and fan-out modes;
their performance remains unmeasured. A proxy demonstration that combines retention, mirroring and
replay after a backend failure remains planned.

Store performance comparisons and qualification of mapped and multiworker paths remain pending
(\S\ref{sec:eval:planned}).
\todom{measure mapped reading and editing, retained-object reuse and fan-out; evaluate the
planned proxy scenario and store variants on the revised build}
